
\documentclass{article}
\usepackage{colortbl}
\usepackage{makecell}
\usepackage{multirow}
\usepackage{supertabular}

\begin{document}

\newcounter{utterance}

\centering \large Interaction Transcript for game `cladder', experiment `full\_v1.5\_default', episode 1743 with Qwen3.5{-}27B.
\vspace{24pt}

{ \footnotesize  \setcounter{utterance}{1}
\setlength{\tabcolsep}{0pt}
\begin{supertabular}{c@{$\;$}|p{.15\linewidth}@{}p{.15\linewidth}p{.15\linewidth}p{.15\linewidth}p{.15\linewidth}p{.15\linewidth}}
    \# & \multicolumn{2}{c}{Player} && \multicolumn{2}{c}{Game Master} \\
    \hline

    \theutterance \stepcounter{utterance}  
    & & & \multicolumn{4}{p{0.6\linewidth}}{
        \cellcolor[rgb]{0.9,0.9,0.9}{
            \makecell[{{p{\linewidth}}}]{
                \texttt{\tiny{[P1$\langle$GM]}}
                \texttt{You are an expert in causal inference. The following question is not a typical commonsense query, but rather a meticulously designed question created by a professor specializing in causal inference, intended to assess the students' mastery of the course content.} \\
\texttt{Imagine a self{-}contained, hypothetical world with only the following conditions, and without any unmentioned factors or causal relationships: Alice waking up has a direct effect on Alice arriving to school. Traffic has a direct effect on Alice arriving to school. For days when Alice wakes up on time, the probability of arriving to school on time is 46\%. For days when Alice wakes up late, the probability of arriving to school on time is 78\%. For days when Alice wakes up late, would it be more likely to see arriving to school on time if Alice had gotten up on time?} \\
\texttt{Start your answer with 'Yes' or 'No', followed by additional reasoning or evidence to support your explanation.} \\
            }
        }
    }
     \\ \\

    \theutterance \stepcounter{utterance}  
    & \multicolumn{4}{p{0.6\linewidth}}{
        \cellcolor[rgb]{0.561,0.737,0.561}{
            \makecell[{{p{\linewidth}}}]{
                \texttt{\tiny{[P1$\rangle$GM]}}
                \texttt{No} \\
\\ 
\texttt{**Reasoning:**} \\
\\ 
\texttt{In this hypothetical world, we are given the conditional probabilities of the outcome (arriving on time) based on the treatment (waking up on time vs. late). Let \$A\$ represent the event "Alice arrives on time" and \$W\$ represent the event "Alice wakes up on time."} \\
\\ 
\texttt{The problem provides the following data:} \\
\texttt{1.  \$P(A | W) = 0.46\$ (Probability of arriving on time given she woke up on time).} \\
\texttt{2.  \$P(A | \textbackslash{}neg W) = 0.78\$ (Probability of arriving on time given she woke up late).} \\
\\ 
\texttt{The question asks: For days when Alice wakes up late (\$\textbackslash{}neg W\$), would it be more likely to see her arriving on time if she had gotten up on time (\$W\$)?} \\
\\ 
\texttt{This is a counterfactual query asking to compare the observed outcome under the current condition (\$\textbackslash{}neg W\$) with the potential outcome under the alternative condition (\$W\$). Specifically, it asks if \$P(A | do(W)) > P(A | \textbackslash{}neg W)\$, or more simply in this context, if the probability of the outcome under the "on time" condition is higher than the probability under the "late" condition.} \\
\\ 
\texttt{Comparing the two probabilities provided:} \\
\texttt{*   Probability if she wakes up on time: **46\%**} \\
\texttt{*   Probability if she wakes up late: **78\%**} \\
\\ 
\texttt{Since \$0.46 < 0.78\$, the probability of arriving on time is actually **lower** if she wakes up on time compared to when she wakes up late. Therefore, it would **not** be more likely to see her arriving on time if she had gotten up on time; in fact, it would be less likely.} \\
\\ 
\texttt{*Note: While this result (waking up late increasing the chance of arriving on time) seems counter{-}intuitive in the real world, the problem explicitly states to imagine a self{-}contained world with only the provided conditions. In this specific causal structure, the data indicates a negative association between waking up on time and arriving on time, possibly due to the unobserved confounding variable "Traffic" acting in a specific way (e.g., perhaps waking up late forces Alice to take a different, faster route, or the "Traffic" variable is perfectly correlated with waking up late in a way that mitigates delays, though the prompt says "without any unmentioned factors," implying we must strictly adhere to the provided probabilities as the ground truth of this world).*} \\
\\ 
\texttt{Based strictly on the numbers provided: \$46\textbackslash{}\% < 78\textbackslash{}\%\$, so the answer is No.} \\
            }
        }
    }
    & & \\ \\

    \theutterance \stepcounter{utterance}  
    & & & \multicolumn{2}{p{0.3\linewidth}}{
        \cellcolor[rgb]{0.9,0.9,0.9}{
            \makecell[{{p{\linewidth}}}]{
                \texttt{\tiny{[GM$|$GM]}}
                \texttt{No} \\
            }
        }
    }
    & & \\ \\

    \theutterance \stepcounter{utterance}  
    & & & \multicolumn{2}{p{0.3\linewidth}}{
        \cellcolor[rgb]{0.9,0.9,0.9}{
            \makecell[{{p{\linewidth}}}]{
                \texttt{\tiny{[GM$|$GM]}}
                \texttt{no} \\
            }
        }
    }
    & & \\ \\

    \theutterance \stepcounter{utterance}  
    & & & \multicolumn{2}{p{0.3\linewidth}}{
        \cellcolor[rgb]{0.9,0.9,0.9}{
            \makecell[{{p{\linewidth}}}]{
                \texttt{\tiny{[GM$|$GM]}}
                \texttt{game\_result = WIN} \\
            }
        }
    }
    & & \\ \\

\end{supertabular}
}

\end{document}
